% TCC — Trabalhos Relacionados
\section{Trabalhos Relacionados}\label{sec:relacionados}

Esta seção posiciona o projeto diante de ferramentas que também testam APIs
REST a partir de contratos. A comparação é bibliográfica e não representa uma
execução conjunta das ferramentas. O objetivo é identificar capacidades e
limites relevantes para o recorte do TCC, sem produzir um ranking entre
soluções com objetivos diferentes.

Os critérios usados são: fonte do modelo de entrada; modo de acesso à
aplicação; estratégia de geração; tratamento de estado e dependências; forma
do oráculo; possibilidade de reproduzir os resultados; e organização ou
priorização dos achados.

\subsection{Ferramentas selecionadas}

\paragraph{Schemathesis.}
Schemathesis deriva testes baseados em propriedades a partir de schemas de
APIs. A geração usa tipos e restrições do contrato e aplica verificações
semânticas às respostas observadas \cite{hatfielddodds2022schemathesis}. A
abordagem é black-box e inclui mecanismos para reduzir e reproduzir casos que
falham. Ela é a ferramenta mais próxima do projeto no uso direto do contrato
para gerar entradas e verificar propriedades, mas seu foco é o teste baseado
em propriedades, não a atribuição de uma prioridade fuzzy aos achados.

\paragraph{RESTler.}
RESTler compila uma especificação Swagger/OpenAPI em uma gramática e infere
dependências produtor--consumidor entre operações. Com isso, gera sequências
stateful que criam o estado necessário antes de exercitar parâmetros
\cite{atlidakis2019restler}. O trabalho também descreve checkers e mecanismos
de agrupamento de bugs. RESTler é uma referência para dependências entre
operações; o fuzzer deste TCC cobre esse aspecto apenas parcialmente e não
pretende oferecer a mesma exploração stateful.

\paragraph{RESTest.}
RESTest realiza teste black-box orientado por OpenAPI e combina geração
aleatória com restrições adicionais. O framework contempla casos nominais e
de falha e associa oráculos aos casos gerados \cite{martinlopez2020restest}.
Essa abordagem ajuda a separar o uso do contrato da implementação de um
oráculo, mas depende do modelo e das restrições disponíveis. O projeto se
aproxima no teste negativo e se diferencia pelo catálogo de mutações e pela
priorização posterior dos achados.

\paragraph{RestTestGen.}
RestTestGen apresenta uma arquitetura extensível para teste black-box, com
parser OpenAPI, provedores de valores, geradores, mutadores, executor e
oráculos \cite{corradini2022resttestgen}. A decomposição é útil para comparar
componentes do pipeline e para discutir reprodução dos artefatos. O trabalho
não define um classificador fuzzy de prioridade; essa é uma parte específica
da proposta deste TCC.

\subsection{Síntese da comparação}

\begin{table}[ht]
\centering
\scalefont{0.82}
\caption{Posicionamento bibliográfico das ferramentas selecionadas}
\label{tab:ferramentas-relacionadas}
\begin{tabular}{|p{2.3cm}|p{2.3cm}|p{1.7cm}|p{2.3cm}|p{2.4cm}|p{2.4cm}|}
\hline
\textbf{Ferramenta} & \textbf{Entrada e geração} & \textbf{Acesso} & \textbf{Estado} & \textbf{Oráculo} & \textbf{Prioridade} \\ \hline
Schemathesis & Schema + property-based & Black-box & Fluxos conforme configuração & Checks e propriedades & Não identificada na fonte-base \\ \hline
RESTler & Gramática + sequências & Black-box & Stateful, produtor--consumidor & Checkers e respostas & Não é foco \\ \hline
RESTest & OpenAPI + restrições & Black-box & Dependente do modelo & Oráculos dos casos & Não é foco \\ \hline
RestTestGen & Framework modular + mutadores & Black-box & Extensível & Oráculos extensíveis & Não identificada \\ \hline
Projeto & OpenAPI + mutações tipadas & Black-box & Parcial & Contrato e detectores & FIS Mamdani pós-execução \\ \hline
\end{tabular}
\end{table}

As quatro ferramentas mostram que “usar OpenAPI” não define uma única técnica.
Elas diferem na maneira de gerar dados, montar sequências, lidar com estado e
interpretar respostas. Os trabalhos de revisão também apontam dificuldades
recorrentes: dependências entre operações, autenticação, dados realistas e
reprodução das execuções. Esses fatores precisam ser controlados ou registrados
para que duas campanhas possam ser comparadas \cite{golmohammadi2023rest,zhang2023openproblems}.

O projeto não reivindica novidade por gerar casos a partir de OpenAPI, testar
entradas negativas ou verificar status e schema. A proposta investigada é mais
específica: combinar mutações estruturadas, evidências reproduzíveis e um
sistema fuzzy que ordena os achados depois da execução. A metodologia avalia
essa combinação no corpus definido; não há comparação experimental suficiente
para afirmar superioridade sobre as ferramentas relacionadas.

\subsection{Critério para eventual comparação experimental}

Schemathesis é a candidata mais próxima caso o orientador considere útil uma
comparação prática. Para ser válida, ela teria de usar os mesmos contratos,
alvos, número de campanhas, condições de reset e orçamento aproximado. As
métricas deveriam permanecer limitadas às definidas no TCC: requisições
concluídas, divergências observadas, famílias deduplicadas, estabilidade e
reprodução dos artefatos. Cobertura interna, qualidade de implementação ou
superioridade de uma ferramenta não seriam inferidas sem um protocolo próprio.

Por ora, RESTler, RESTest e RestTestGen cumprem papel de referências
arquiteturais e metodológicas. EvoMaster e Dredd permanecem registrados no
material de apoio, mas não entram nesta comparação até que sua utilidade e a
fonte primária adequada sejam confirmadas.
